% ---------------------------------------------------------------------
% 00-introduction.tex
% The introduction does the following, in this order:
%   1. states the problem and the main theorem, in full;
%   2. places the result against the literature, with citations;
%   3. explains the method and the obstruction it overcomes;
%   4. outlines the structure of the paper.
% ---------------------------------------------------------------------

\section{Introduction}
\label{sec:introduction}

Replace this paragraph with a statement of the problem. Fix the
notation and the ambient setting, then state the main result as a
displayed theorem so the reader can find it immediately. The house
convention is to state the principal theorem once here and again,
with full proof, in \cref{sec:main-result}; the introduction copy may
restate it in a slightly more discursive form.

The literature discussion should be specific. Attribute each earlier
result to its authors with a citation, and say precisely which
hypothesis is dropped, which range is extended, or which construction
is replaced by the present argument. For example, the classical
treatment of this circle of ideas is due to earlier authors
\cite{example-classical}, and the closest modern statement is
\cite{example-modern}.

A second paragraph should describe the proof strategy. Name the key
new ingredient, say where it enters, and identify the point at which
the argument would fail without it. Avoid promising more than the
paper delivers.

Finally, outline the remaining sections:
\begin{itemize}
  \item \Cref{sec:preliminaries} collects definitions and standing
        conventions.
  \item \Cref{sec:main-result} proves the main theorem.
  \item \Cref{sec:conclusion} records consequences and open problems.
  \item \Cref{app:sample-appendix} contains routine verifications.
\end{itemize}
